% ===================== 3. KHỞI TẠO LỜI GIẢI =====================
\section{Chiến lược khởi tạo lời giải ban đầu}
\label{sec:init}

Greedy tạo nghiệm khởi đầu \emph{hạt giống} dùng chung cho cả ba tầng. Mục tiêu:
nhanh, hợp lệ, và đủ tốt để các tầng sau chỉ cần cải thiện.

\subsection{Sắp xếp ưu tiên}
Thứ tự xử lý quyết định chất lượng Greedy, theo hai tiêu chí:
\begin{itemize}
  \item \textbf{Thứ tự lớp:} sắp theo \emph{thời lượng $t$ tăng dần}, rồi \emph{sĩ
  số $s$ tăng dần} --- ưu tiên lớp \emph{ngắn và ít sinh viên} trước. Lớp ngắn linh
  hoạt, dễ lấp vào các khe thời gian nhỏ, nên xếp sớm giúp giữ lại các khối thời
  gian dài liền mạch cho những lớp khó hơn về sau.
  \item \textbf{Thứ tự phòng (best-fit):} duyệt phòng theo \emph{sức chứa tăng dần}
  --- luôn thử phòng \emph{vừa đủ} chứa trước, để dành các phòng lớn cho những lớp
  đông sinh viên thực sự cần đến chúng, tránh lãng phí tài nguyên khan hiếm.
\end{itemize}

\subsection{Greedy: First-feasible insertion}
Với mỗi lớp theo thứ tự ưu tiên, duyệt tiết bắt đầu hợp lệ $s_0\in\mathcal{S}(t_i)$
(vòng ngoài) rồi phòng theo sức chứa tăng dần (vòng trong); kiểm tra giáo viên rồi
phòng có rảnh cả khối $t_i$ tiết không, và lấy cặp $(\text{tiết},\text{phòng})$
hợp lệ \emph{đầu tiên} để gán luôn. Chiến lược ``first-feasible'' (thay vì tìm vị
trí tối ưu) giữ chi phí ở mức thấp nhất --- đây là lý do Greedy chạy trong
\emph{micro-giây} ngay cả với $N=1000$.

\begin{figure}[H]
\centering
\begin{tikzpicture}[
  font=\scriptsize,
  cls/.style={draw, thick, fill=cfast!15, draw=cfast!70, minimum height=0.55cm, inner sep=2pt},
  room/.style={draw, thick, minimum width=1.5cm, minimum height=0.7cm, fill=black!6},
  chosen/.style={room, fill=cexact!18, draw=cexact, very thick},
  dimmed/.style={room, fill=black!3, draw=gray!55, text=gray},
  hd/.style={font=\small\bfseries, anchor=west},
  lcap/.style={font=\tiny, align=center},
  arr/.style={->, thick},
]
% ---------- Bước 1: hàng đợi đã sắp ----------
\node[hd] at (0,2.25) {\textcolor{accent}{1.} Sắp ưu tiên: $t\uparrow$, rồi $s\uparrow$};
\node[cls, minimum width=0.5cm] (q1) at (0.55,0.7) {};
\node[cls, minimum width=0.5cm, right=4pt of q1] (q2) {};
\node[cls, minimum width=1.0cm, right=4pt of q2] (q3) {};
\node[cls, minimum width=1.6cm, right=4pt of q3] (q4) {};
\node[lcap, below=1pt of q1] {$t{=}1$\\$s{=}20$};
\node[lcap, below=1pt of q2] {$t{=}1$\\$s{=}40$};
\node[lcap, below=1pt of q3] {$t{=}2$\\$s{=}30$};
\node[lcap, below=1pt of q4] {$t{=}3$\\$s{=}60$};
\draw[arr] ($(q1.north west)+(0,0.3)$) -- ($(q4.north east)+(0,0.3)$)
  node[midway, above, font=\tiny\itshape] {thứ tự xử lý (ngắn \& nhỏ trước)};

% ---------- Bước 2: best-fit phòng ----------
\node[hd] at (0,-1.15) {\textcolor{accent}{2.} Best-fit: phòng vừa đủ chứa trước};
\node[cls, minimum width=1.0cm] (cc) at (0.7,-2.2) {$s{=}30$};
\node[chosen, right=1.1cm of cc] (ra) {$c{=}40$};
\node[room, right=5pt of ra] (rb) {$c{=}80$};
\node[dimmed, right=5pt of rb] (rc) {$c{=}200$};
\draw[arr, cexact, very thick] (cc) -- (ra) node[midway, above, font=\tiny] {chọn};
\node[lcap, text=cexact!60!black, below=1pt of ra] {vừa đủ ($\ge 30$)};
\node[lcap, text=gray, below=1pt of rc] {để dành lớp đông};
\node[lcap, right=0.5cm of rc, text width=2.6cm, align=left] {\itshape Phòng sắp theo $c\uparrow$};
\end{tikzpicture}
\caption{Chiến lược khởi tạo Greedy. \textbf{(1)} Lớp được xử lý theo thứ tự
\emph{thời lượng tăng, sĩ số tăng} (ngắn \& ít SV trước). \textbf{(2)} Với mỗi lớp,
quét tiết bắt đầu (vòng ngoài) rồi phòng theo \emph{sức chứa tăng dần} và chọn phòng
\emph{vừa đủ} chứa đầu tiên (best-fit, first-feasible) --- để dành phòng lớn cho lớp
đông. Quá trình lặp lại đến khi không xếp thêm được lớp nào.}
\label{fig:greedy-init}
\end{figure}

\subsection{Lặp gán đến điểm dừng}
\label{sec:repair}
Greedy không gán một lần rồi dừng: nó \emph{lặp lại} các lượt quét cho đến khi một
lượt không xếp thêm được lớp nào (điểm dừng). Cơ chế lặp này gắn liền với Local
Search ở Mục~\ref{sec:ls}: sau mỗi lần Ejection Chain tái sắp xếp và giải phóng
chỗ, một lượt Greedy mới có thể tận dụng để nhồi thêm lớp.

\begin{remark}
Greedy không cố tối ưu --- nó cố \emph{nhanh và hợp lệ}. Chính vì còn dư địa cải
thiện (đặc biệt trên dữ liệu nghẽn), ba tầng phía sau mới có việc để làm; và vì cả
ba cùng xuất phát từ một mốc Greedy, mọi so sánh ở Mục~\ref{sec:results} đều quy về
cùng một điểm tham chiếu.
\end{remark}
